% Generated draft: requires mathematical and source review.
\documentclass[11pt]{article}
\usepackage[margin=25mm]{geometry}
\usepackage{fontspec}
\setmainfont{Latin Modern Roman}
\usepackage{amsmath,amssymb,mathtools,unicode-math}
\setmathfont{Latin Modern Math}
\usepackage{enumitem,needspace,xurl,hyperref}
\setlength{\parindent}{0pt}
\setlength{\parskip}{5pt}
\setlength{\emergencystretch}{3em}
\newcommand{\N}{\mathbb{N}}
\newcommand{\Z}{\mathbb{Z}}
\newcommand{\Q}{\mathbb{Q}}
\newcommand{\R}{\mathbb{R}}
\newcommand{\C}{\mathbb{C}}
\newcommand{\F}{\mathbb{F}}
\newcommand{\A}{\mathbb{A}}
\newcommand{\PP}{\mathbb{P}}
\newcommand{\E}{\mathbb{E}}
\newcommand{\eps}{\varepsilon}
\newcommand{\dd}{\,\mathrm{d}}
\begin{document}
\section*{UnsolvedMath ID: 30001320}
\textbf{Source problem code:} catalogue entry
\subsection{Definitions and mathematical statement}
For a sequence of finite box graphs \(B_n=(V_n,E_n)\) of increasing size in a fixed dimension \(d\), free boundary conditions mean that only edges in \(E_n\), with both endpoints inside the box, are available: no periodic edges or boundary identifications are added. Under independent bond percolation with parameter \(p\), each edge is open independently with probability \(p\). An open cluster is a connected component of the graph consisting of all vertices and the open edges, including isolated vertices. Define the largest cluster size by
\[
M_n=\max\bigl\{|C|:C\text{ is an open cluster in }B_n\bigr\}.
\]
Here size is interpreted as the number of vertices; the supplied record does not explicitly specify this convention.

The supplied problem asks: in high dimension, at criticality and with free boundary conditions, determine the asymptotic order of \(M_n\) as the box size tends to infinity, and characterize its limiting behavior. A precise version requires specifying the box graphs, the permitted dimensions, the critical parameter, and the intended notion of asymptotic order and convergence. For example, a distributional formulation would ask for deterministic normalizations \(a_n>0\), possibly centerings \(b_n\), and a limiting law \(\mu\) such that
\[
\frac{M_n-b_n}{a_n}\xrightarrow{\mathrm d}\mu.
\]
This displays a possible meaning of limiting behavior; the source does not explicitly require this formulation, assert that such a limit exists, or specify whether it should be nondegenerate.

The supplied statement concerns a box with free boundaries, whereas the supplied literature assessment refers to a ``free-boundary high-dimensional torus.'' These descriptions do not specify the same boundary geometry without further explanation. Consequently, this draft preserves the box formulation while recording the unresolved discrepancy; it does not substitute a periodic-boundary problem.
\subsection{Short English statement}
At criticality, determine how the largest open cluster grows in a high-dimensional box with free boundaries and describe its limiting behavior. The graph, dimension range, meaning of criticality, and type of limit remain unspecified. The literature assessment also introduces an unresolved box-versus-torus discrepancy.
\subsection{Sources}
\begin{itemize}
\item Local UnsolvedMath registry, permanent ID \(30001320\), source problem code \texttt{catalogue entry}, titled ``Largest Critical Percolation Clusters with Free Boundaries.'' The original, clean, and main statement fields agree. The registry labels the statement ``exact'' and supplies a manual-verification assertion; those labels are supplied provenance, not independent mathematical verification of this draft.
\item Supplied report citation: \emph{Scaling Limits in Models of Statistical Mechanics} (2010), Oberwolfach Reports, DOI \texttt{10.4171/owr/2009/38}, pp.\ 2115--2178. Supplied source URL and literature-source link, labeled ``original Oberwolfach report 4081'': \url{https://doi.org/10.4171/owr/2009/38}. Attribution is to the report's workshop open-problem discussions; no individual proposer is supplied. The proposed year is recorded as 2009, while the citation gives 2010.
\item Literature assessment dated exactly \texttt{2026-08-21}: the supplied curation labels the problem open and states that the report identifies a free-boundary high-dimensional torus percolation asymptotic as open, with later results cited there not supplying the missing boundary-condition theorem. This is supplied dated curation, not an independently verified literature assessment. The source itself requests independent verification of time-sensitive status and cited claims.
\end{itemize}
\end{document}
